\section{相关研究与方法定位}
\subsection{资本分散与风险分散}
Markowitz 的组合选择理论以预期收益和协方差刻画资产之间的配置关系\cite{markowitz1959portfolio}。这一框架为比较资本分散与风险分散提供了共同的收益风险语言。

均值与协方差构成资产配置的基本统计描述。资本等权容易实现，但当各资产波动不同，等资本并不对应等风险。相反，若将更多资本配置给低波动工具，组合的风险贡献可能更加均衡，而现金或债券的资本占比也可能明显提高。两种分散概念回答不同问题，不能只选择其中一种作为所有模型的评价标准。

风险预算研究以资产对组合风险的边际影响为基础。Qian 的风险平价讨论和后续风险预算文献为这种配置逻辑提供了系统表达\cite{qian2005risk,roncalli2013introduction}。Maillard 等对等风险贡献组合的性质展开研究\cite{maillard2010properties}。这些工作说明，相关性参与决定风险贡献，单纯按波动倒数分配只有在特定结构下才与等风险贡献一致。

\subsection{宽松风险平价的基本思想}
Gambeta 与 Kwon 在宽松风险平价中引入目标收益，研究偏离风险平价与组合表现之间的权衡\cite{gambeta2020risk}。Ardia 等从绩效与风险贡献的平衡出发，讨论风险型组合的扩展\cite{ardia2018beyond}。本文沿用收益与风险预算可以权衡的研究思路，但采用可行参考收益目标和两阶段凸问题，不将本文目标函数与上述文献的具体模型等同。

将风险预算作为完全不可违反的等式，可能限制投资者对收益和其他实施条件的表达。宽松处理允许在保留风险预算解释的同时接受一定偏离。这里的“宽松”不意味着放松信息时序或核算标准，而是将最终组合与风险预算参考之间的差异纳入可控目标。

本文采用两阶段表达。第一阶段只构造风险预算参考，第二阶段在多头预算集合上求解参考跟踪、组合方差和收益短缺的加权目标。参考点承担锚定作用，最终权重承担实施作用。两者均通过凸问题求得，但它们的最优性条件并不相同。将参考权重与最终权重分开，有助于解释为什么最终组合可能持有某些资产为零。

\subsection{凸优化与可审计性}
凸优化提供的是给定输入下的求解结构。若目标函数为凸函数、可行域为凸集合，局部最优点可以具有全局最优意义；严格凸性还可用于讨论解的唯一性。Boyd 与 Vandenberghe 的教材系统讨论了这些条件及最优性判据\cite{boyd2004convex}。本文依据这些规则检查实际程序，而非仅凭目标函数中出现二次项便认定整个问题为凸。

对数风险预算问题中的正二次项与负对数项可分别验证凸性。最终配置中的跟踪项是平方范数，风险项是半正定二次型，收益短缺是仿射函数正部的平方。预算等式与非负约束保持凸可行域。归一化尺度必须在优化开始前由历史输入固定，不能随着待优化权重变化而在分母中再次引入非线性比例。

\subsection{协方差收缩与收益估计}
Chopra 与 Ziemba 研究了均值、方差和协方差估计误差对组合选择的影响\cite{chopra1993effect}。Laloux 等进一步指出，有限样本相关矩阵中的特征值可能受到显著噪声影响\cite{laloux1999noise}。因此，优化问题的数值可解性与输入的统计可靠性需要分别检验。

协方差估计需要同时处理方差与相关性。当有效观察相对资产数量有限，样本中的细小变化可能改变特征值和参考权重。Ledoit--Wolf 方法将经验协方差与结构化目标组合，通过样本估计收缩强度，以改善矩阵估计条件\cite{ledoit2004well}。本文只采用这一指定估计方式，不展示其他协方差配置的绩效。

收益估计采用 20 个交易日半衰期的指数加权均值。这里的指数衰减作用于预测输入，不意味着投资者预测未来恰好持续 20 日，也不意味着组合每天交易。主模型仍按周度日历决策。在同一历史窗口内，近期收益对均值估计影响更大；而协方差的收缩机制控制的是矩阵结构，两者承担不同任务。

\subsection{层次配置与简单对照}
DeMiguel 等将等权规则与多种组合优化方法作样本外比较，说明估计误差可能抵消优化带来的分散收益\cite{demiguel2009optimal}。Jagannathan 与 Ma 对非负权重约束的研究也表明，约束可能具有缓和估计误差的作用\cite{jagannathan2003risk}。这些结果支持保留简单对照和明确的资本约束，但不意味着任何受约束模型都能获得样本外优势。

HRP Benchmark 以相关性距离组织资产，再通过递归划分分配资本。相关文献强调层次结构在分散配置中的作用\cite{lopezdeprado2016building}。HERC Benchmark 的文献背景进一步涉及层次结构中的风险贡献\cite{raffinot2018hierarchical}。本文仓库实现采用排序后二分和簇波动比例，是明确的简化递归规则；本文不将其与原文全部算法步骤等同。

Equal Weight 保留简单的等资本参照，60/40 Benchmark 保留固定权益与债券预算参照。它们可以帮助辨别高收益是否仅来自更多资本暴露于高波动资产，以及低换手是否来自较少的信号更新。对照不需要在每个指标上输给主模型才能有价值，差异本身构成对研究问题的回答。

\subsection{尾部风险与实施研究}
CVaR 将注意力集中于损失分布的尾部。Rockafellar 与 Uryasev 给出了适合优化的表示\cite{rockafellar2000optimization}。本文将其用于事后风险统计，并未在指定 Global RRP 中启用 CVaR 约束。因而后文的尾部损失较低只能作为路径特征，不能归因为一个未实际存在的优化限制。

交易成本同样需要区分目标函数与收益扣减。本文最终求解器不包含换手惩罚，周度交易的成本在日收益中直接扣除。这样的结果能够回答约定费率下的净表现，但尚不能回答真实资金规模下的容量。文献与实现共同提示，投资者需要分别评价配置方法、估计质量和执行条件，不能将历史净值曲线视作对所有环节的同时验证。

Rockafellar 与 Uryasev 对一般损失分布的讨论说明，有限样本和离散损失下的尾部风险定义仍需保持一致\cite{rockafellar2002conditional}。Almgren 与 Chriss 将交易执行中的价格风险、暂时性冲击和永久性冲击纳入分析\cite{almgren2001optimal}，为区分本文的固定费率扣减与实际执行成本提供了依据。

\subsection{国内实践与研究有效性}
国内风险平价研究关注资产波动差异、工具范围和配置实施条件。华泰证券关于国内实践的研究讨论了资产风险与宏观风险暴露的联系\cite{htsc2024practice}，低利率环境研究则将多资产配置与绝对收益目标相联系\cite{htsc2025lowrate}。这些研究提供了应用背景，本文仍以无杠杆 ETF 组合为研究对象，不将其他策略的收益结果作为本模型的验证。

历史策略比较还涉及多重检验与规格选择。White 讨论了重复使用数据进行模型选择所产生的数据窥探问题\cite{white2000reality}；Bailey 等提出回测过拟合概率框架\cite{bailey2014probability}，Bailey 与 López de Prado 的调整夏普方法进一步考虑选择偏差与非正态性\cite{bailey2017pbo}。本文据此区分逐期信息时序合规与研究规格的历史选择，不将一倍标准误候选集合解释为已经完成多重检验校正，也不报告当前配置未计算的过拟合概率或调整夏普数值。
